\usepackage{fontspec}
\IfFontExistsTF{Times New Roman}
  {\setmainfont{Times New Roman}}
  {\setmainfont{Liberation Serif}}
\IfFontExistsTF{Arial}
  {\setsansfont{Arial}}
  {\setsansfont{Liberation Sans}}
% Türkçe babel'in etkin "=" kısaltması paket seçeneklerini bozmasın.
\usepackage[shorthands=off,turkish]{babel}
\usepackage{unicode-math}
\IfFontExistsTF{Cambria Math}
  {\setmathfont{Cambria Math}}
  {\setmathfont{TeX Gyre Termes Math}}

% Tüm kimlik bilgilerini yalnızca bu dosyada güncelleyin.
\newcommand{\TezBasligi}{Enerji Depolama Sistemleri İçin Geliştirilmiş T-Tipi Dönüştürücü Mimarisi ve Düşük Güç Performansını İyileştiren Anahtarlama ve Sin\textsuperscript{3} Tabanlı Kompanzasyon Yöntemleri}
\newcommand{\TezBasligiPDF}{Enerji Depolama Sistemleri İçin Geliştirilmiş T-Tipi Dönüştürücü Mimarisi ve Düşük Güç Performansını İyileştiren Anahtarlama ve Sin3 Tabanlı Kompanzasyon Yöntemleri}
\newcommand{\TezBasligiKapak}{ENERJİ DEPOLAMA SİSTEMLERİ İÇİN GELİŞTİRİLMİŞ T-TİPİ DÖNÜŞTÜRÜCÜ MİMARİSİ\\VE DÜŞÜK GÜÇ PERFORMANSINI İYİLEŞTİREN ANAHTARLAMA VE\\Sin\textsuperscript{3} TABANLI KOMPANZASYON YÖNTEMLERİ}
\newcommand{\TezYazari}{Ümit PEKPARLAK}
\newcommand{\OgrenciNo}{[DOKTORA ÖĞRENCİ NUMARASI]}
\newcommand{\AnabilimDali}{[ANABİLİM DALI]}
\newcommand{\BilimDali}{[BİLİM DALI / PROGRAM]}
\newcommand{\Danisman}{Prof. Dr. Uğur ARİFOĞLU}
\newcommand{\OrtakDanisman}{}
\newcommand{\TeslimTarihi}{HAZİRAN 2026}
\newcommand{\SavunmaTarihi}{[SAVUNMA TARİHİ]}


\usepackage{geometry}
\geometry{
  a4paper,
  twoside,
  inner=40mm,
  outer=25mm,
  top=25mm,
  bottom=25mm,
  footskip=12.5mm
}

\usepackage{setspace}
\onehalfspacing
\setlength{\parindent}{0pt}
\setlength{\parskip}{6pt}
\setlength{\emergencystretch}{2em}

\usepackage{fancyhdr}
\fancypagestyle{sau}{%
  \fancyhf{}
  % Metin alanı aynalı olsa da numara kâğıdın fiziksel merkezinde kalır.
  \fancyhfoffset[LO,RE]{15mm}
  \fancyfoot[C]{\normalfont\bfseries\fontsize{12}{12}\selectfont\thepage}
  \renewcommand{\headrulewidth}{0pt}
  \renewcommand{\footrulewidth}{0pt}
}
% Birinci derece başlıkların ilk sayfasında numara gösterilmez.
\fancypagestyle{plain}{%
  \fancyhf{}
  \renewcommand{\headrulewidth}{0pt}
  \renewcommand{\footrulewidth}{0pt}
}
% Çift taraflı baskıdaki boş sayfalar tamamen boş kalır.
\makeatletter
\renewcommand{\cleardoublepage}{%
  \clearpage
  \if@twoside
    \ifodd\c@page\else
      \hbox{}\thispagestyle{empty}\newpage
      \if@twocolumn\hbox{}\newpage\fi
    \fi
  \fi
}
\makeatother

\usepackage{graphicx}
\usepackage{float}
% Yeni hattın kendi şekilleri figures/ altında; mevcut kanıt görselleri
% tez-latex/figures deposundan okunur (çoğaltılmaz).
\graphicspath{{figures/}{../tez-latex/figures/}}
\usepackage[table]{xcolor}

% ---------------------------------------------------------------------------
% EKSİK ÇALIŞMA barı: tamamlanması istenen MATLAB benzetimi, veri veya
% analiz işini sonucun geleceği yerde görünür kılar. Nihai teslimden önce
% bütün barlar ya kapatılmış ya da Sınırlılıklar bölümüne taşınmış olmalıdır.
% Kullanım: \eksikistek{Kısa başlık}{Ne istiyoruz — hangi koşu, hangi veri.}
% ---------------------------------------------------------------------------
\newcommand{\eksikistek}[2]{%
  \par\medskip
  \noindent\fcolorbox{orange!85!black}{orange!12}{%
    \begin{minipage}{0.94\linewidth}
      \small
      \colorbox{orange!85!black}{\strut\textcolor{white}{\textbf{EKSİK ÇALIŞMA --- BUNU İSTİYORUZ}}}\ \textbf{#1}\par\smallskip
      #2
    \end{minipage}%
  }%
  \par\medskip
}

% ---------------------------------------------------------------------------
% ÜMİT ABİ İÇİN NOT kutusu (K-28): tez sahibine bilgi notu; tez metninin
% parçası değildir. Teslimden önce aşağıdaki satır \umitnotlarifalse yapılır;
% bütün notlar tek noktadan PDF'ten kalkar (metin akışı değişmez).
% Kullanım: \umitnot{Kısa başlık}{Not metni.}
% ---------------------------------------------------------------------------
\newif\ifumitnotlari
\umitnotlaritrue
% Kutu sayfa sonunda bölünebilir (tcolorbox, breakable); paket yalnız notlar
% açıkken yüklenir, notsuz çıktı etkilenmez.
\ifumitnotlari
  \RequirePackage[breakable,skins]{tcolorbox}
  \newtcolorbox{umitnotkutusu}[1]{breakable, enhanced jigsaw, colback=blue!6, colframe=blue!60!black,
    boxrule=0.6pt, arc=1pt, left=4pt, right=4pt, top=3pt, bottom=3pt, fontupper=\small,
    title={\textbf{ÜMİT ABİ İÇİN NOT --- TESLİMDE ÇIKARILACAK}\quad #1}, fonttitle=\small}
\fi
\newcommand{\umitnot}[2]{%
  \ifumitnotlari
  \par\medskip
  \begin{umitnotkutusu}{#1}
    #2
  \end{umitnotkutusu}
  \par\medskip
  \fi
}

\usepackage{booktabs}
\usepackage{longtable}
\usepackage{tabularx}
\usepackage{array}
% unicode-math matematik sembollerini sağlar; amssymb ile \eth çakışmasını önleriz.
\usepackage{amsmath}
\usepackage{siunitx}
\sisetup{output-decimal-marker={,}}

\usepackage{caption}
\DeclareCaptionFormat{sauhang}{#1#2#3}
\captionsetup{
  font={normalfont,stretch=1},
  labelfont=bf,
  labelsep=period,
  format=hang,
  justification=justified,
  singlelinecheck=true
}
\captionsetup[table]{position=top,skip=6pt}
\captionsetup[figure]{position=bottom,skip=6pt}
\usepackage{subcaption}
\usepackage{tikz}
\usetikzlibrary{arrows.meta,positioning,fit,calc,decorations.pathmorphing}
\numberwithin{equation}{chapter}

\usepackage{enumitem}
\setlist{itemsep=6pt,topsep=6pt,parsep=0pt,partopsep=0pt}

% Sayısal atıfları metinde görünme sırasına göre sıkıştırarak gösterir.
\usepackage[sort&compress,numbers]{natbib}

\usepackage[absolute,overlay]{textpos}
\textblockorigin{0mm}{0mm}

\usepackage[hidelinks]{hyperref}
\hypersetup{
  pdftitle={\TezBasligiPDF},
  pdfauthor={\TezYazari}
}

\usepackage{titlesec}
% SAÜ FBE: 1. derece başlık 12 punto, kalın, büyük harf, sola yaslı.
\titleformat{\chapter}[hang]
  {\normalfont\bfseries\fontsize{12}{12}\selectfont}
  {\thechapter.}{0.6em}{}
\titlespacing*{\chapter}{0pt}{72pt}{18pt}

% 2. derece: önce 18 nk, sonra 12 nk.
\titleformat{\section}
  {\normalfont\bfseries\fontsize{12}{12}\selectfont}
  {\thesection.}{0.6em}{}
\titlespacing*{\section}{0pt}{18pt}{12pt}

% 3. ve 4. derece: önce 12 nk, sonra 6 nk.
\titleformat{\subsection}
  {\normalfont\bfseries\fontsize{12}{12}\selectfont}
  {\thesubsection.}{0.6em}{}
\titlespacing*{\subsection}{0pt}{12pt}{6pt}
\titleformat{\subsubsection}
  {\normalfont\bfseries\fontsize{12}{12}\selectfont}
  {\thesubsubsection.}{0.6em}{}
\titlespacing*{\subsubsection}{0pt}{12pt}{6pt}

\setcounter{secnumdepth}{3}
\setcounter{tocdepth}{3}

% Resmî Word şablonundaki İçindekiler stilleri:
% tek satır, paragraf öncesi/sonrası 0 nk ve tüm düzeylerde noktalı kılavuz.
\usepackage{tocloft}
\setlength{\cftparskip}{0pt}
\setlength{\cftbeforechapskip}{0pt}
\setlength{\cftbeforesecskip}{0pt}
\setlength{\cftbeforesubsecskip}{0pt}
\setlength{\cftbeforesubsubsecskip}{0pt}
\renewcommand{\cftchapdotsep}{\cftdotsep}
\renewcommand{\cftchapfont}{\normalfont\bfseries}
\renewcommand{\cftchappagefont}{\normalfont\bfseries}
\renewcommand{\cfttoctitlefont}{\normalfont\bfseries\fontsize{12}{12}\selectfont}
\renewcommand{\cftlottitlefont}{\normalfont\bfseries\fontsize{12}{12}\selectfont}
\renewcommand{\cftloftitlefont}{\normalfont\bfseries\fontsize{12}{12}\selectfont}
\setlength{\cftbeforetoctitleskip}{72pt}
\setlength{\cftaftertoctitleskip}{18pt}
\setlength{\cftbeforelottitleskip}{72pt}
\setlength{\cftafterlottitleskip}{18pt}
\setlength{\cftbeforeloftitleskip}{72pt}
\setlength{\cftafterloftitleskip}{18pt}
\setlength{\cftchapindent}{0pt}
\setlength{\cftchapnumwidth}{22.51mm}
\setlength{\cftsecindent}{4mm}
\setlength{\cftsubsecindent}{8.5mm}
\setlength{\cftsubsubsecindent}{12.7mm}

% Numarasız birinci derece ön/arka bölüm başlığı.
\newcommand{\saufrontchapter}[2][]{%
  \cleardoublepage
  \chapter*{#2}
  \thispagestyle{plain}
  \if\relax\detokenize{#1}\relax\else
    \addcontentsline{toc}{chapter}{#1}
  \fi
}

% Biçim sistemi, 04.08.2022 tarih ve 621-9 sayılı Senato kararı eki
% SAÜ FBE tez yazım kılavuzu, Word şablonu ve teslim kontrol formundan
% distile edilmiştir. Kontrol formu ile kılavuzun sayfa numarası mesafesi
% çeliştiğinden Word şablonundaki 1,25 cm değeri uygulanmıştır.
